\begin{problem}[33届物理竞赛复赛]{56}
    如图，上、下两个平凸透光柱面的半径分别为 $R_{1}$ 、 $R_{2}$ ，且两柱面外切；其剖面（平面）分别平行于各自的轴线，且相互平行；各自过切点的母线相互垂直。取两柱面切点O为直角坐标系O-XYZ的原点，下侧柱面过切点O的母线为X轴，上侧柱面过切点O的母线为Y轴。一束在真空中波长为 $\lambda$ 的可见光沿Z轴负方向傍轴入射，分别从上、下柱面反射回来的光线会发生干涉；借助于光学读数显微镜，逆着Z轴方向，可观测到原点附近上方柱面上的干涉条纹在X-Y平面的投影。 $R_{1}$ 和 $R_{2}$ 远大于傍轴光线干涉区域所对应的两柱面间最大间隙。空气折射率为 $n_{0}=1.00$ 。试推导第k级亮纹在X-Y平面的投影的曲线方程。

    已知：a. 在两种均匀、各向同性的介质的分界面两侧，折射率较大（小）的介质为光密（疏）介质；光线在光密（疏）介质的表面反射时，反射波存在（不存在）半波损失。任何情形下，折射波不存在半波损失。伴随半波损失将产生大小为 $\pi$ 的相位突变。b. $\sin x \approx x$ ，当 $|x| << 1$ 。
\end{problem}